% Bagian: sets-functions-relations
% Bab: size-of-sets
% Subbagian: enumerations

\documentclass[../../../include/open-logic-section]{subfiles}

\begin{document}

\olfileid[id]{sfr}{siz}{enm}

\olsection{Enumerasi dan Himpunan \usetoken{S}{enumerable}}

\begin{editorial}
  Bagian ini membahas enumerasi himpunan dengan mendefinisikannya sebagai
  surjeksi dari $\PosInt$. Uraiannya disampaikan secara perlahan bagi pembaca
  dengan sedikit latar belakang matematika. Versi alternatif yang lebih
  ringkas diberikan dalam \olref[enm-alt]{sec}; versi itu mendefinisikan
  enumerasi secara berbeda, yakni sebagai bijeksi dengan $\Nat$ (atau suatu
  segmen awal).
\end{editorial}

\begin{explain}
Kita telah memberi contoh himpunan dengan mencantumkan !!{element}snya.
Sekarang mari kita bahas secara lebih umum bagaimana dan kapan kita dapat
mencantumkan !!{element}s suatu himpunan, sekalipun himpunan itu tak hingga.
\end{explain}

\begin{defn}[Enumerasi, secara informal]
Secara informal, suatu \emph{enumerasi} himpunan~$A$ adalah daftar (yang
mungkin tak hingga) dari !!{element}s~$A$ sedemikian sehingga setiap
!!{element} dari $A$ muncul pada suatu posisi hingga dalam daftar. Jika $A$
mempunyai suatu enumerasi, maka $A$ disebut \emph{!!{enumerable}}.
\end{defn}

\begin{explain}
Beberapa hal mengenai enumerasi:
\begin{enumerate}
\item Kita hanya menganggap sebagai enumerasi daftar yang mempunyai awal dan
  yang setiap !!{element} selain anggota pertama mempunyai tepat satu
  !!{element} tepat sebelumnya. Dengan kata lain, hanya ada sejumlah hingga
  anggota di antara !!{element} pertama dalam daftar dan sebarang
  !!{element} lainnya. Secara khusus, ini berarti bahwa setiap !!{element}
  dalam suatu enumerasi mempunyai posisi hingga: !!{element} pertama
  berposisi~$1$, anggota kedua berposisi~$2$, dan seterusnya.
\item Kita dapat mempunyai enumerasi yang berbeda untuk himpunan~$A$ yang sama
  karena urutan kemunculan !!{element}snya berbeda: $4$, $1$, $25$, $16$,
  dan~$9$ mengenumerasi (himpunan) lima bilangan kuadrat pertama sama baiknya
  dengan $1$, $4$, $9$, $16$, dan~$25$.
\item Enumerasi yang memuat pengulangan tetap merupakan enumerasi: $1$, $1$,
  $2$, $2$, $3$, $3$,~\dots{} mengenumerasi himpunan yang sama seperti $1$,
  $2$, $3$,~\dots{}.
\item Urutan dan pengulangan \emph{memang} penting ketika kita menentukan suatu
  enumerasi: kita dapat mengenumerasi bilangan bulat positif mulai dengan
  $1$, $2$, $3$, $1$, \dots{}, tetapi polanya lebih mudah dilihat jika
  dienumerasi dengan cara baku sebagai $1$, $2$, $3$, $4$,~\dots
\item Enumerasi harus mempunyai awal: \dots, $3$, $2$, $1$ bukanlah enumerasi
  bilangan bulat positif karena daftar itu tidak mempunyai !!{element}
  pertama. Untuk melihat bagaimana hal ini mengikuti definisi informal,
  tanyakan kepada diri Anda, ``pada posisi berapa dalam daftar bilangan 76
  muncul?''
\item Daftar berikut bukan enumerasi bilangan bulat positif: $1$, $3$, $5$,
  \dots, $2$, $4$, $6$, \dots\@ Masalahnya ialah bilangan genap menempati
  posisi $\infty + 1$, $\infty + 2$, $\infty + 3$, alih-alih posisi hingga.
\item Himpunan kosong terhitung: himpunan itu dienumerasi oleh daftar kosong!{}
\end{enumerate}
\end{explain}

\begin{prop}
Jika $A$ mempunyai suatu enumerasi, maka himpunan itu mempunyai enumerasi tanpa
  pengulangan.
\end{prop}

\begin{proof}
  Andaikan $A$ mempunyai enumerasi $x_1$, $x_2$, \dots{} yang setiap $x_i$-nya
  merupakan !!{element} dari~$A$. Kita dapat menghilangkan pengulangan dari
  enumerasi dengan membuang !!{element}s yang berulang. Sebagai contoh, kita
  dapat mengubah enumerasi itu menjadi enumerasi baru: kita mencantumkan $x_i$
  jika ia adalah !!a{element} dari~$A$ yang tidak terdapat di antara $x_1$,
  \dots, $x_{i-1}$, atau membuang $x_i$ dari daftar jika ia telah muncul di
  antara $x_1$, \dots,~$x_{i-1}$.
\end{proof}

Argumen terakhir menunjukkan bahwa untuk memahami dengan baik enumerasi dan
himpunan yang !!{enumerable}, serta membuktikan berbagai hal tentangnya, kita
memerlukan definisi yang lebih saksama. Definisi berikut memberikannya.

\begin{defn}[Enumerasi, secara formal]
Suatu \emph{enumerasi} dari himpunan $A \neq \emptyset$ adalah sebarang
fungsi !!{surjective} $f \colon \PosInt \to A$.
\end{defn}

\begin{explain}
Mari kita yakinkan diri bahwa definisi formal dan definisi informal yang
menggunakan daftar yang mungkin tak hingga itu ekuivalen. Pertama, setiap
fungsi !!{surjective} dari $\PosInt$ ke suatu himpunan~$A$ mengenumerasi~$A$.
Fungsi semacam itu menentukan enumerasi sebagaimana didefinisikan secara
informal di atas: daftar $f(1)$, $f(2)$, $f(3)$, \dots. Karena $f$
!!{surjective}, setiap !!{element} dari~$A$ pasti merupakan nilai~$f(n)$ untuk
suatu~$n \in \PosInt$. Karena itu, setiap !!{element} dari $A$ muncul pada
suatu posisi hingga dalam daftar. Karena fungsi tersebut mungkin tidak
!!{injective}, daftar itu mungkin mengandung pengulangan, tetapi hal itu dapat
diterima (sebagaimana telah disebutkan di atas).

Sebaliknya, jika diberikan suatu daftar yang mengenumerasi semua !!{element}s
dari~$A$, kita dapat mendefinisikan fungsi !!a{surjective} $f\colon \PosInt
\to A$ dengan menetapkan $f(n)$ sebagai !!{element} ke-$n$ dalam daftar, atau
!!{element} terakhir dalam daftar jika tidak ada !!{element} ke-$n$. Satu-
satunya kasus yang tidak menghasilkan fungsi !!a{surjective} ialah ketika $A$
kosong, sehingga daftarnya pun kosong. Jadi, setiap daftar tak kosong
menentukan fungsi !!a{surjective} $f\colon \PosInt \to A$.
\end{explain}

\begin{defn}
  \ollabel{defn:enumerable}
  Suatu himpunan~$A$ adalah !!{enumerable} jika dan hanya jika himpunan itu kosong atau
  mempunyai suatu enumerasi.
\end{defn}

\begin{ex}
Fungsi yang mengenumerasi bilangan bulat positif ($\PosInt$) cukup berupa
fungsi identitas yang diberikan oleh $f(n) = n$. Fungsi yang mengenumerasi
bilangan asli $\Nat$ adalah fungsi $g(n) = n - 1$.
\end{ex}

\begin{ex}
Fungsi $f\colon \PosInt \to \PosInt$ dan $g \colon \PosInt \to \PosInt$ yang
diberikan oleh
\begin{align*}
f(n) & = 2n \text{ dan}\\
g(n) & = 2n - 1
\end{align*}
masing-masing mengenumerasi bilangan bulat positif genap dan bilangan bulat
positif ganjil. Namun, tidak satu pun merupakan enumerasi dari $\PosInt$,
karena tidak satu pun !!{surjective}.
\end{ex}

\begin{prob}
Definisikan suatu enumerasi dari bilangan kuadrat positif $1$, $4$, $9$,
$16$, \dots
\end{prob}

\begin{ex}
Fungsi $f(n) = (-1)^{n} \lceil \frac{(n-1)}{2}\rceil$ (dengan $\lceil x
\rceil$ menyatakan fungsi \emph{pembulatan ke atas}, yang membulatkan $x$ ke
bilangan bulat terdekat di atasnya) mengenumerasi himpunan bilangan
bulat~$\Int$. Perhatikan bagaimana $f$ menghasilkan nilai-nilai $\Int$ dengan
``melompat'' bolak-balik antara bilangan bulat positif dan negatif:
\[
\begin{array}{c c c c c c c c}
f(1) & f(2) & f(3) & f(4) & f(5) & f(6) & f(7) & \dots \\ \\
- \lceil \tfrac{0}{2} \rceil & \lceil \tfrac{1}{2}\rceil & - \lceil \tfrac{2}{2} \rceil & \lceil \tfrac{3}{2} \rceil & - \lceil \tfrac{4}{2} \rceil  & \lceil \tfrac{5}{2}
\rceil & - \lceil \tfrac{6}{2} \rceil & \dots \\ \\
0 & 1 & -1 & 2 & -2 & 3 & \dots
\end{array}
\]
Kita juga dapat memandang $f$ sebagai fungsi yang didefinisikan menurut kasus
sebagai berikut:
\[
f(n) = \begin{cases}
  0 & \text{jika $n = 1$}\\
  n/2 & \text{jika $n$ genap}\\
  -(n-1)/2 & \text{jika $n$ ganjil dan $>1$}
  \end{cases}
\]
\end{ex}

\begin{prob}
  Tunjukkan bahwa jika $A$ dan $B$ adalah !!{enumerable}, maka $A \cup B$ juga
  demikian. Untuk melakukannya, andaikan ada fungsi !!{surjective} $f\colon
  \PosInt \to A$ dan $g\colon \PosInt \to B$, lalu definisikan fungsi
  !!a{surjective}~$h\colon \PosInt \to A \cup B$ dan buktikan bahwa fungsi itu
  !!{surjective}. Tinjau pula kasus ketika $A$ atau~$B = \emptyset$.
\end{prob}

\begin{prob}
  Tunjukkan bahwa jika $B \subseteq A$ dan $A$ adalah !!{enumerable}, maka $B$
  juga demikian. Untuk melakukannya, andaikan ada fungsi !!a{surjective}
  $f\colon \PosInt \to A$. Definisikan fungsi !!a{surjective}~$g\colon
  \PosInt \to B$ dan buktikan bahwa fungsi itu !!{surjective}. Apa yang
  terjadi jika $B = \emptyset$?
\end{prob}

\begin{prob}
  Tunjukkan dengan induksi pada $n$ bahwa jika $A_1$, $A_2$, \dots, $A_n$
  semuanya !!{enumerable}, maka $A_1 \cup \dots \cup A_n$ juga demikian. Anda
  boleh menggunakan fakta bahwa jika dua himpunan $A$ dan~$B$ adalah
  !!{enumerable}, maka $A \cup B$ juga demikian.
\end{prob}

Meskipun ketika mencantumkan !!{element}s suatu himpunan mungkin lebih alami
untuk mulai menghitung dari !!{element} ke-$1$, matematikawan suka menggunakan
bilangan asli~$\Nat$ untuk menghitung. Mereka berbicara tentang !!{element}s
ke-$0$, ke-$1$, ke-$2$, dan seterusnya dalam suatu daftar. Sejalan dengan itu,
kita dapat mendefinisikan enumerasi sebagai fungsi !!a{surjective} dari $\Nat$
ke~$A$. Tentu saja, kedua definisi itu ekuivalen.

\begin{prop}\ollabel{prop:enum-shift}
  Terdapat !!a{surjection} $f\colon \PosInt \to A$ jika dan hanya jika
  terdapat !!a{surjection} $g\colon \Nat \to A$.
\end{prop}

\begin{proof}
  Jika diberikan !!a{surjection} $f\colon \PosInt \to A$, kita dapat
  mendefinisikan $g(n) = f(n+1)$ untuk semua $n \in \Nat$. Mudah dilihat bahwa
  $g\colon \Nat \to A$ adalah !!{surjective}. Sebaliknya, jika diberikan
  !!a{surjection} $g\colon \Nat \to A$, definisikan $f(n) = g(n-1)$.
\end{proof}

Ini memberi kita hasil berikut:

\begin{cor}\ollabel{cor:enum-nat}
Suatu himpunan $A$ adalah !!{enumerable} jika dan hanya jika himpunan itu kosong atau
terdapat fungsi !!a{surjective} $f\colon \Nat \to A$.
\end{cor}

Di atas kita telah membahas bahwa daftar !!{element}s suatu himpunan~$A$
dapat diubah menjadi daftar tanpa pengulangan. Hal ini juga benar untuk
enumerasi, tetapi sedikit lebih sulit dirumuskan dan dibuktikan secara ketat.
Setiap fungsi $f\colon \PosInt \to A$ harus terdefinisi untuk semua $n \in
\PosInt$. Jika hanya ada sejumlah hingga !!{element}s dalam~$A$, jelas kita
tidak mungkin mempunyai fungsi yang terdefinisi pada !!{element}s $\PosInt$
yang tak terhingga banyaknya, mengambil semua !!{element}s $A$ sebagai nilai,
tetapi tidak pernah mengambil nilai yang sama dua kali. Dalam kasus itu,
yakni ketika daftar tanpa pengulangan tersebut hingga, kita harus memilih
domain lain untuk~$f$, yang hanya mempunyai sejumlah hingga !!{element}s.
Tidak adanya pengulangan berarti bahwa $f$ harus !!{injective}. Karena fungsi itu
juga !!{surjective}, kita mencari !!a{bijection} antara suatu himpunan hingga
$\{1, \dots, n\}$ atau $\PosInt$ dengan~$A$.

\begin{prop}\ollabel{prop:enum-bij}
Jika $f\colon \PosInt \to A$ adalah !!{surjective} (yakni, suatu enumerasi
dari~$A$), terdapat !!a{bijection} $g\colon Z \to A$ dengan $Z$ berupa
$\PosInt$ atau $\{1, \dots, n\}$ untuk suatu~$n \in \PosInt$.
\end{prop}

\begin{proof}
  Kita mendefinisikan fungsi $g$ secara rekursif: Tetapkan $g(1) = f(1)$. Jika
  $g(i)$ telah didefinisikan, tetapkan $g(i+1)$ sebagai nilai pertama di antara
  $f(1)$, $f(2)$, \dots{} yang belum terdapat di antara $g(1)$, \dots, $g(i)$,
  jika nilai semacam itu ada. Jika $A$ hanya mempunyai $n$ !!{element}s, maka
  $g(1)$, \dots, $g(n)$ semuanya terdefinisi, sehingga kita telah
  mendefinisikan fungsi $g\colon \{1, \dots, n\} \to A$. Jika $A$ mempunyai
  tak hingga banyak !!{element}s, maka untuk setiap $i$ harus ada !!a{element}
  dari~$A$ dalam enumerasi $f(1)$, $f(2)$, \dots, yang belum terdapat di antara
  $g(1)$, \dots, $g(i)$. Dalam kasus ini kita telah mendefinisikan fungsi
  $g\colon \PosInt \to A$.

  Fungsi $g$ adalah !!{surjective}, karena setiap anggota dari~$A$ terdapat di
  antara $f(1)$, $f(2)$, \dots{} (karena $f$ !!{surjective}), dan pada akhirnya
  akan menjadi nilai~$g(i)$ untuk suatu~$i$. Fungsi itu juga !!{injective},
  karena jika terdapat $j < i$ sedemikian sehingga $g(j) = g(i)$, maka $g(i)$
  telah terdapat di antara $g(1)$, \dots, $g(i-1)$, bertentangan dengan cara
  kita mendefinisikan~$g$.
\end{proof}

\begin{cor}\ollabel{cor:enum-nat-bij}
Suatu himpunan $A$ adalah !!{enumerable} jika dan hanya jika himpunan itu kosong atau
terdapat !!a{bijection} $f\colon N \to A$, dengan $N = \Nat$ atau $N = \{0,
\dots, n\}$ untuk suatu $n \in \Nat$.
\end{cor}

\begin{proof}
$A$ adalah !!{enumerable} jika dan hanya jika $A$ kosong atau terdapat
!!a{surjective} $f\colon \PosInt \to A$. Menurut \olref{prop:enum-bij}, hal
yang terakhir berlaku jika dan hanya jika terdapat fungsi !!a{bijective}
~$f\colon Z \to A$ dengan $Z = \PosInt$ atau $Z = \{1, \dots, n\}$ untuk suatu
$n \in \PosInt$. Dengan argumen yang sama seperti dalam bukti
\olref{prop:enum-shift}, hal itu pada gilirannya berlaku jika dan hanya jika
terdapat !!a{bijection} $g\colon N \to A$, dengan $N = \Nat$ atau $N = \{0,
\dots, n-1\}$.
\end{proof}

\begin{prob}
  Menurut \olref[sfr][siz][enm]{defn:enumerable}, suatu himpunan $A$ adalah
  terhitung jika dan hanya jika $A = \emptyset$ atau terdapat !!a{surjective}
  $f\colon \PosInt \to A$. Kita juga dapat mendefinisikan ``himpunan
  !!{enumerable}'' secara saksama sebagai berikut: suatu himpunan terhitung
  jika dan hanya jika terdapat fungsi !!a{injective} $g\colon A \to \PosInt$.
  Tunjukkan bahwa kedua definisi tersebut ekuivalen, yakni tunjukkan bahwa
  terdapat fungsi !!a{injective} $g\colon A \to \PosInt$ jika dan hanya jika
  $A = \emptyset$ atau terdapat !!a{surjective} $f\colon \PosInt \to A$.
\end{prob}

\end{document}
